% Vietnamese translation for OI Public Library
% Source-File: IOI中国国家候选队论文/2007/day1/9.王欣上《浅谈基于分层思想的网络流算法》.ppt
\documentclass[11pt]{article}
\usepackage{oipl-vn}

\newcommand{\Oh}{\mathcal{O}}
\newcommand{\level}{\operatorname{level}}

\title{Bản trình chiếu: Thuật toán luồng mạng dựa trên phân tầng}
\author{Wang Xinshang}
\date{2007}

\begin{document}
\maketitle

\origtitle{Slide companion for layered network-flow algorithms}
\sourcepath{IOI China candidate team papers / 2007 / day 1 / Wang Xinshang.ppt}

\section{Phạm vi bản dịch}

Tệp PowerPoint đi kèm bài viết đã được dịch từ PDF về các thuật toán luồng
mạng dựa trên đồ thị phân tầng. Bản ghi chú này giữ lại các khái niệm và mạch
thuật toán chính của slide.

\section{Đồ thị dư và mức}

Với một luồng hiện tại, đồ thị dư biểu diễn lượng luồng còn có thể đẩy thêm
theo cạnh thuận và lượng luồng có thể hoàn lại theo cạnh ngược. Chạy BFS từ
nguồn trên đồ thị dư cho ta mức của mỗi đỉnh:
\[
  \level(u)=\text{độ dài đường ngắn nhất từ nguồn tới }u.
\]
Đồ thị phân tầng chỉ giữ các cạnh dư \((u,v)\) thỏa
\[
  \level(v)=\level(u)+1.
\]

\section{Luồng chặn}

Trong một pha, thuật toán tìm luồng chặn trên đồ thị phân tầng, tức đẩy luồng
cho tới khi không còn đường từ nguồn tới đích chỉ đi qua các cạnh tăng mức.
Sau pha đó, độ dài đường tăng luồng ngắn nhất trong đồ thị dư tăng lên, nên
số pha bị chặn bởi số đỉnh.

\section{Các thuật toán}

Slide nhắc ba tên chính: MPLA, Dinic và MPM. Dinic là thuật toán quen thuộc
nhất trong thi lập trình: mỗi pha dựng mức bằng BFS, sau đó DFS nhiều lần
trên đồ thị phân tầng để tìm luồng chặn. Con trỏ hiện tại giúp bỏ qua cạnh đã
hết khả năng, làm cài đặt nhanh hơn.

MPM và các biến thể khác cũng dựa trên cùng quan điểm phân tầng, nhưng chọn
cách tính luồng chặn khác. Khi cần một cài đặt ổn định trong thi, Dinic
thường là lựa chọn đầu tiên.

\section{Ứng dụng}

Tư tưởng phân tầng đặc biệt hữu ích khi mô hình mạng có nhiều cạnh nhưng các
đường tăng luồng ngắn mang nhiều thông tin. Các bài ghép cặp, phủ, chọn đối
tượng dưới ràng buộc và nhiều mô hình cắt cực tiểu có thể dùng Dinic như một
thành phần trung tâm sau bước dựng mạng.

\end{document}
